\section{Thiết kế ứng dụng Flutter}
Ứng dụng Flutter được thiết kế để người dùng theo dõi trạng thái căn hộ và điều khiển thiết bị từ điện thoại. Giao diện chính hiển thị các thông số nhiệt độ, độ ẩm, CO2, ánh sáng, TDS và pH. Mỗi thông số có màu trạng thái để người dùng nhanh chóng nhận biết bình thường, cảnh báo hoặc vượt ngưỡng.

Ứng dụng có các màn hình chức năng chính:
\begin{itemize}
    \item \textbf{Màn hình tổng quan:} hiển thị giá trị mới nhất và trạng thái môi trường.
    \item \textbf{Màn hình biểu đồ:} hiển thị lịch sử thay đổi của từng thông số theo thời gian.
    \item \textbf{Màn hình điều khiển:} bật/tắt quạt, bật/tắt đèn, tăng/giảm nhiệt độ điều hòa.
    \item \textbf{Màn hình cấu hình:} thiết lập ngưỡng cảnh báo cho CO2, nhiệt độ, độ ẩm, ánh sáng, TDS và pH.
\end{itemize}

Ứng dụng không giao tiếp trực tiếp với node mà kết nối với ThingsBoard. Cách thiết kế này giúp người dùng có thể truy cập từ xa, đồng thời giảm độ phức tạp của node và gateway. Mọi dữ liệu và lệnh đều được đồng bộ qua nền tảng IoT, bảo đảm hệ thống có thể mở rộng khi bổ sung thêm node trong tương lai.

\subsection{Màn hình đăng nhập}
Ảnh giao diện đăng nhập cho thấy ứng dụng có tên hiển thị \textit{Smart Farm}, biểu tượng lá cây và hai trường nhập email, password. Cách đặt tên này phù hợp với hướng hệ thống chăm sóc cây thông minh, tuy nhiên trong phạm vi khóa luận có thể hiểu app là giao diện điều khiển node cây trồng trong nhà và chất lượng không khí. Màn hình đăng nhập giúp giới hạn quyền truy cập, tránh người ngoài điều khiển relay hoặc điều hòa.

\begin{figure}[h!]
    \centering
    \begin{tikzpicture}[every node/.style={font=\small}]
        \draw[rounded corners=18pt, fill=green!8] (0,0) rectangle (5,8);
        \node[fill=green!45, circle, minimum size=0.9cm] at (2.5,6.3) {};
        \node[font=\bfseries\Large, text=green!50!black] at (2.5,5.6) {Smart Farm};
        \draw[rounded corners=10pt, fill=white] (0.55,2.2) rectangle (4.45,4.8);
        \draw[rounded corners=5pt] (0.85,4.05) rectangle (4.15,4.55);
        \node[anchor=west] at (1.1,4.3) {Email};
        \draw[rounded corners=5pt] (0.85,3.25) rectangle (4.15,3.75);
        \node[anchor=west] at (1.1,3.5) {Password};
        \draw[rounded corners=5pt, fill=green!55] (0.85,2.65) rectangle (4.15,3.1);
        \node[text=white] at (2.5,2.88) {Đăng nhập};
        \node at (2.5,2.35) {Quên mật khẩu?};
    \end{tikzpicture}
    \caption{Minh họa màn hình đăng nhập của ứng dụng}
    \label{fig:app_login}
\end{figure}

\subsection{Màn hình Sensor}
Màn hình Sensor hiển thị sáu thẻ dữ liệu gồm TDS, pH, AS, CO2, nhiệt độ và độ ẩm. Dữ liệu trong ảnh lần lượt là TDS 698.86 ppm, pH 7.1, ánh sáng 0.66 lux, CO2 2655 ppm, nhiệt độ 30.9 $^\circ$C và độ ẩm 65\%. Đây là nhóm thông số phản ánh đồng thời tình trạng cây và chất lượng không khí. CO2 ở mức 2655 ppm cho thấy phòng có thể đang thiếu thông gió; ánh sáng 0.66 lux cho thấy vị trí cây gần như thiếu sáng; nhiệt độ 30.9 $^\circ$C tương đối cao đối với không gian sinh hoạt.

\begin{figure}[h!]
    \centering
    \begin{tikzpicture}[every node/.style={font=\small}]
        \draw[rounded corners=14pt, fill=green!5] (0,0) rectangle (6,9);
        \node[anchor=west, font=\Large] at (0.3,8.45) {Node};
        \node[anchor=east, font=\bfseries] at (4.8,8.45) {Manual};
        \draw[rounded corners=8pt, fill=gray!25] (5.0,8.2) rectangle (5.7,8.65);
        \node[text=green!50!black] at (1.5,7.7) {Sensor};
        \node at (4.5,7.7) {Actuator};
        \draw[green!50!black, thick] (1.1,7.45) -- (1.9,7.45);
        \foreach \x/\y/\t/\v in {
            0.4/5.6/TDS/698.86 ppm,
            3.2/5.6/PH/7.1,
            0.4/3.65/AS/0.66 lux,
            3.2/3.65/CO2/2655.0 ppm,
            0.4/1.7/Nhiệt độ/30.9 $^\circ$C,
            3.2/1.7/Độ ẩm/65.0 \%
        }{
            \draw[rounded corners=8pt, fill=green!22, draw=green!22] (\x,\y) rectangle ++(2.4,1.65);
            \node[text=green!55!black, font=\Large] at (\x+1.2,\y+1.05) {$\bullet$};
            \node at (\x+1.2,\y+0.65) {\t};
            \node[font=\bfseries] at (\x+1.2,\y+0.32) {\v};
        }
    \end{tikzpicture}
    \caption{Minh họa màn hình Sensor hiển thị dữ liệu node}
    \label{fig:app_sensor}
\end{figure}

\subsection{Màn hình Actuator}
Màn hình Actuator hiển thị ba thiết bị gồm RL1, RL2 và điều hòa. Trong ảnh, cả ba thiết bị đều ở trạng thái OFF. Bố cục dạng thẻ đồng nhất với màn hình Sensor, giúp người dùng dễ chuyển đổi giữa giám sát và điều khiển. Khi người dùng nhấn vào một thẻ actuator, app gửi lệnh tương ứng đến ThingsBoard.

\begin{figure}[h!]
    \centering
    \begin{tikzpicture}[every node/.style={font=\small}]
        \draw[rounded corners=14pt, fill=green!5] (0,0) rectangle (6,8);
        \node[anchor=west, font=\Large] at (0.3,7.45) {Node};
        \node[anchor=east, font=\bfseries] at (4.8,7.45) {Manual};
        \draw[rounded corners=8pt, fill=gray!25] (5.0,7.2) rectangle (5.7,7.65);
        \node at (1.5,6.7) {Sensor};
        \node[text=green!50!black] at (4.5,6.7) {Actuator};
        \draw[green!50!black, thick] (4.1,6.45) -- (4.9,6.45);
        \foreach \x/\y/\t in {
            0.4/4.8/RL1,
            3.2/4.8/RL2,
            0.4/2.85/Điều hòa
        }{
            \draw[rounded corners=8pt, fill=green!22, draw=green!22] (\x,\y) rectangle ++(2.4,1.65);
            \node[text=green!55!black, font=\Large] at (\x+1.2,\y+1.05) {$\clubsuit$};
            \node at (\x+1.2,\y+0.65) {\t};
            \node[font=\bfseries] at (\x+1.2,\y+0.32) {OFF};
        }
    \end{tikzpicture}
    \caption{Minh họa màn hình Actuator điều khiển relay và điều hòa}
    \label{fig:app_actuator}
\end{figure}

\subsection{Ánh xạ dữ liệu từ ThingsBoard sang giao diện}
App Flutter cần ánh xạ telemetry từ ThingsBoard sang các thành phần giao diện. Mỗi thẻ Sensor tương ứng với một khóa telemetry. Nếu khóa không tồn tại hoặc dữ liệu quá cũ, app nên hiển thị trạng thái mất dữ liệu thay vì giữ nguyên giá trị cũ mà không cảnh báo. Điều này đặc biệt quan trọng với CO2 và nhiệt độ vì người dùng có thể đưa ra quyết định dựa trên dữ liệu hiển thị.

\begin{longtable}{|p{3.2cm}|p{3.5cm}|p{6cm}|}
\hline
\textbf{Thẻ giao diện} & \textbf{Telemetry} & \textbf{Cách hiển thị} \\
\hline
TDS & \texttt{tds} & Hiển thị ppm, cảnh báo khi vượt ngưỡng cây. \\
\hline
PH & \texttt{ph} & Hiển thị một chữ số thập phân hoặc hai chữ số tùy hiệu chuẩn. \\
\hline
AS & \texttt{lux} & Hiển thị lux, cảnh báo thiếu sáng. \\
\hline
CO2 & \texttt{co2} & Hiển thị ppm, cảnh báo khi phòng thiếu thông gió. \\
\hline
Nhiệt độ & \texttt{temperature} & Hiển thị $^\circ$C, phục vụ điều khiển điều hòa. \\
\hline
Độ ẩm & \texttt{humidity} & Hiển thị \%, cảnh báo khô hoặc ẩm cao. \\
\hline
\end{longtable}

\subsection{Bảo mật và quyền điều khiển}
Vì app có khả năng điều khiển thiết bị thật, bảo mật không thể bỏ qua. Màn hình đăng nhập giúp xác thực người dùng trước khi truy cập dữ liệu. Token ThingsBoard không nên hard-code trực tiếp trong mã nguồn public. App cần lưu token đăng nhập trong bộ nhớ an toàn của thiết bị và xóa token khi người dùng đăng xuất.

Ngoài xác thực, hệ thống cần phân quyền. Người dùng thông thường có thể xem dữ liệu và điều khiển thủ công; tài khoản quản trị có thể thay đổi ngưỡng, cấu hình node hoặc thêm thiết bị. Phân quyền rõ ràng giúp hạn chế rủi ro khi nhiều người cùng sử dụng hệ thống trong một căn hộ.
